\begin{tcolorbox}[title=Problem 10, breakable]
    If $F$ is a field and $\rho$ is a ring homomorphism, is $\rho(F)$
    also a field? If yes, prove it? If no, give a counterexample.
\end{tcolorbox} 

\begin{proof}
    Consider the mapping $f : \mathbb{R} \longrightarrow \mathbb{R}$ defined 
    \[f(x) = 0 \text{ for all } x \in \mathbb{R}.\]
\end{proof}

\begin{tcolorbox}[title=Problem 12, breakable]
    If $I$ is an ideal of a commutative ring $R$, then show that two cosets of $I$
    (say $I + a$ and $I + b$) are either equal or disjoint. (That is the 
    set of all translates of $I$ \emph{partition} $R$.) If you've previously 
    encountered the idea of equivalence relation, rephrase this result in terms of that concept.
\end{tcolorbox} 

\begin{proof}
    Suppose $I+a \cap I+b \neq \emptyset$. Then there exists $x \in R$ such that
    \[
        x \in I+a \quad \text{and} \quad x \in I+b.
    \]
    Thus there exist $i_1, i_2 \in I$ such that
    \[
        x = a + i_1 = b + i_2.
    \]
    Rearranging shows
    \[
        a - b = i_2 - i_1.
    \]
    Since $I$ is an ideal, it is closed under subtraction, so $i_2 - i_1 \in I$.
    Thus $a - b \in I$, which implies $a \in b + I$, so $I+a = I+b$.

    In terms of an equivalence relation, define $a \sim b$ if $a-b \in I$.
    Then the equivalence classes under $\sim$ are exactly the cosets of $I$.
\end{proof}

\begin{tcolorbox}[title=Problem 16, breakable]
    Let $R$ be a finite commutative ring, with $n$ elements.
    Suppose that $I$ is an ideal of $R$ with $m$ elements.
    Prove that $m$ divides $n$.
\end{tcolorbox} 

\begin{proof}
    Notice $R / I$ divides $R$ into partitions 
        of equivalent cardinality because each coset has the same cardinality as $I$.
    By Problem 16 these partitions are disjoint and their union is $R$.
    Thus $|I| x = |R|$ for some $x = |R/I|$.
\end{proof}